% eksperymentalne tłumaczenie na włoski

%

\section{Łatwe} % lang-pl
\section{Problemi facili} % lang-it


Hartshorne \cite[s. 122]{hartshorne2000} napisze, że Kartezjusz około 1637 pokazał, że: % lang-pl
Hartshorne \cite[p. 122]{hartshorne2000} scrive che Cartesio mostrò intorno al 1637 che: % lang-it

\begin{proposition}
    Niech $P_1 = (x_1, y_1)$, $P_2 = (x_2, y_2)$, \ldots, $P_n = (x_n, y_n)$ będą punktami na rzeczywistej płaszczyźnie kartezjańskiej, na której dane są również punkty $(0, 0)$ oraz $(1, 0)$. % lang-pl
    
    Wtedy następujące warunki są równoważne: % lang-pl
    Siano $P_1 = (x_1, y_1)$, $P_2 = (x_2, y_2)$, \ldots, $P_n = (x_n, y_n)$ punti del piano cartesiano reale, sul quale sono inoltre assegnati i punti $(0,0)$ e $(1,0)$. % lang-it
    
    Allora le seguenti condizioni sono equivalenti: % lang-it
    \begin{itemize}
        \item punkt $P = (x, y)$ można wykreślić przy pomocy cyrkla i linijki; % lang-pl
        \item liczby $x, y$ można otrzymać z $a_1, a_2, \ldots, a_n, b_1, b_2, \ldots, b_n$ przy pomocy dodawania, odejmowania, mnożenia, dzielenia lub rozwiązując skończoną liczbę liniowych lub kwadratowych równań, w których wyciągamy pierwiastki kwadratowe jedynie z dodatnich liczb. % lang-pl
        \item il punto $P=(x,y)$ può essere costruito con riga e compasso; % lang-it
        \item i numeri $x,y$ possono essere ottenuti da $a_1, a_2, \ldots, a_n, b_1, b_2, \ldots, b_n$ mediante addizione, sottrazione, moltiplicazione, divisione oppure risolvendo un numero finito di equazioni lineari o quadratiche, nelle quali si estraggono radici quadrate soltanto di numeri positivi. % lang-it
    \end{itemize}
\end{proposition}

% TODO: Hartshorne s. 103
\begin{geoconstruction}
    Dany jest kąt. % lang-pl
    
    Znaleźć jego dwusieczną. % lang-pl
    È dato un angolo. % lang-it
    
    Costruirne la bisettrice. % lang-it
\index{dwusieczna}%
\end{geoconstruction}

Hartshorne \cite[s. 23]{hartshorne2000}. % lang-pl

Hartshorne \cite[p. 23]{hartshorne2000}. % lang-it
% TODO: Zetel s. 21 to samo, ale wierzchołek jest niedostępny

\begin{problem}
    Środek odcinka. % lang-pl
    
    Il punto medio di un segmento. % lang-it
\end{problem}

Hartshorne \cite[s. 23]{hartshorne2000}. % lang-pl

Hartshorne \cite[p. 23]{hartshorne2000}. % lang-it

\begin{problem}
    Prosta prostopadła do prostej, przechodząca przez punkt. % lang-pl
    
    La retta perpendicolare a una retta assegnata passante per un punto. % lang-it
    \index{prosta!prostopadła}
\end{problem}

Hartshorne \cite[s. 23, 24]{hartshorne2000}. % lang-pl

Hartshorne \cite[p. 23, 24]{hartshorne2000}. % lang-it

\begin{problem}
    Prosta równoległa do prostej, przechodząca przez punkt. % lang-pl
    
    La retta parallela a una retta assegnata passante per un punto. % lang-it
    \index{prosta!równoległa}
\end{problem}

Hartshorne \cite[s. 24]{hartshorne2000}. % lang-pl
Hartshorne \cite[p. 24]{hartshorne2000}. % lang-it
\begin{problem} % lang-it
    Il centro di un cerchio assegnato. % lang-it
\end{problem} % lang-it


\begin{problem} % lang-pl
    Środek danego okręgu. % lang-pl
\end{problem} % lang-pl
Hartshorne \cite[p. 24]{hartshorne2000}. % lang-it

Hartshorne \cite[s. 24]{hartshorne2000}. % lang-pl
Porównaj z zadaniem Napoleona \ref{napoleon_problem}. % lang-pl
Confrontare con il problema di Napoleone \ref{napoleon_problem}. % lang-it
\index{zadanie!Napoleona}

\begin{problem}
    Styczna do okręgu przez punkt poza nim. % lang-pl
    
    La tangente a un cerchio passante per un punto esterno. % lang-it
    \index{styczna}
\end{problem}

Hartshorne \cite[s. 24]{hartshorne2000}. % lang-pl
Hartshorne \cite[p. 24]{hartshorne2000}. % lang-it
\begin{problem} % lang-it
    Sono dati due punti $A$, $B$ e una retta $k$. % lang-it
    Trovare un cerchio tangente alla retta e passante per entrambi i punti. % lang-it
\end{problem} % lang-it


Triangolo inscritto: 13 passi, circoscritto: 7 passi. % lang-it
\begin{problem}
    Dane są dwa punkty $A$, $B$ i prosta $k$. % lang-pl
    
    Znaleźć okrąg styczny do prostej, który przechodzi przez obydwa punkty. % lang-pl
    È dato un triangolo con un lato distinto. % lang-it
    
    Inscrivere in esso un quadrato in modo che uno dei suoi lati giaccia sul lato distinto del triangolo. % lang-it
\end{problem}

Trójkąt wpisany 13, opisany 7. % lang-pl
Eves \cite[p. 162]{eves1_1972} % lang-it
\begin{problem} % lang-it
    Sono dati due triangoli, uno contenuto nell'altro. % lang-it
    Inscrivere nel triangolo esterno un terzo triangolo circoscritto a quello interno. % lang-it
\end{problem} % lang-it


Eves \cite[p. 167]{eves1_1972} % lang-it
\begin{problem}
    Dany jest trójkąt z wyróżnionym bokiem. % lang-pl
    
    Wpisać w niego kwadrat tak, by jeden z jego czterech boków leżał na wyróżnionym boku trójkąta. % lang-pl
    Sono dati una retta, un segmento e un punto esterno alla retta. % lang-it
    
    Costruire un triangolo isoscele con vertice nel punto assegnato e base, situata sulla retta, uguale al segmento dato. % lang-it
\end{problem}

Eves \cite[s. 162]{eves1_1972} % lang-pl
Hartshorne \cite[p. 25]{hartshorne2000}. % lang-it
\begin{problem} % lang-it
    Sono dati una retta passante per $B$ e un punto $A$ esterno ad essa. % lang-it
    Trovare un cerchio passante per entrambi i punti e tangente alla retta. % lang-it
\end{problem} % lang-it


Hartshorne \cite[p. 25]{hartshorne2000}. % lang-it
\begin{problem}
    Dane są dwa trójkąty, jeden zawarty w drugim. % lang-pl
    
    Wpisać w zewnętrzny trójkąt trzeci trójkąt tak, by był opisany na wewnętrznym trójkącie. % lang-pl
    Tre cerchi che si intersecano a coppie ortogonalmente. % lang-it
\end{problem}

Eves \cite[s. 167]{eves1_1972} % lang-pl

Problema di Hartshorne \cite[p. 25]{hartshorne2000}; una soluzione di Łukasz Rajkowski in $\Delta_{24}^{12}$ utilizza un fascio di cerchi. % lang-it

\begin{problem}
    Dana jest prosta, odcinek i punkt poza prostą. % lang-pl
    
    Wykreślić trójkąt równoramienny, z wierzchołkiem w punkie i z podstawą równą odcinkowi na prostej. % lang-pl
\end{problem} % lang-pl
Hartshorne \cite[s. 25]{hartshorne2000}. % lang-pl
\begin{problem} % lang-pl
    Dana jest prosta przez punkt $B$, punkt $A$ poza prostą. % lang-pl
    Znaleźć okrąg, który przechodzi przez obydwa punkty, styczny do prostej. % lang-pl
\end{problem} % lang-pl
Hartshorne \cite[s. 25]{hartshorne2000}. % lang-pl
\begin{problem} % lang-pl
    Trzy okręgi przecinające się parami pod kątem prostym. % lang-pl
\end{problem} % lang-pl
Problem Hartshorne'a \cite[s. 25]{hartshorne2000}, rozwiązanie Łukasza Rajkowskiego $\Delta_{24}^{12}$ z wykorzystaniem pęku okręgów. % lang-pl
\begin{problem} % lang-pl
    Podzielić odcinek na trzy równe części. % lang-pl
    Dividere un segmento in tre parti uguali. % lang-it
    \index{trysekcja!odcinka}
\end{problem}

Hartshorne \cite[s. 25]{hartshorne2000}. % lang-pl

Hartshorne \cite[p. 25]{hartshorne2000}. % lang-it

% TODO: Neugebauer s. 67
\begin{problem}
    Okrąg styczny do prostej, przechodzący przez dwa punkty poza nią. % lang-pl
    
    Un cerchio tangente a una retta e passante per due punti esterni ad essa. % lang-it
\end{problem}

Hartshorne \cite[s. 43]{hartshorne2000}. % lang-pl
Hartshorne \cite[p. 43]{hartshorne2000}. % lang-it
\begin{problem} % lang-it
    Un cerchio tangente a due rette e passante per un punto esterno ad esse. % lang-it
\end{problem} % lang-it


Hartshorne \cite[p. 44]{hartshorne2000}. % lang-it
\begin{problem}
    Okrąg styczny do dwóch prostych, przechodzący przez punkt poza nimi. % lang-pl
    
    Un rettangolo avente la stessa area di un triangolo assegnato. % lang-it
    È assegnato uno dei lati del rettangolo. % lang-it
\end{problem}

Hartshorne \cite[s. 44]{hartshorne2000}. % lang-pl
Hartshorne \cite[p. 43]{hartshorne2000}. % lang-it
\begin{problem} % lang-it
    Un quadrato avente la stessa area di un rettangolo assegnato. % lang-it
\end{problem} % lang-it


Hartshorne \cite[p. 43]{hartshorne2000}. % lang-it
\begin{problem}
    Prostokąt o polu takim, jak dany trójkąt. % lang-pl
    
    Dany jest jeden bok prostokąta. % lang-pl
    Dividere un triangolo mediante una retta passante per un punto assegnato in due parti di uguale area. % lang-it
\end{problem}

Hartshorne \cite[s. 43]{hartshorne2000}. % lang-pl

Hartshorne \cite[p. 44]{hartshorne2000}. % lang-it

\begin{problem}
    Kwadrat o polu takim, jak dany prostokąt. % lang-pl
    
\end{problem} % lang-pl
Hartshorne \cite[s. 43]{hartshorne2000}. % lang-pl
\begin{problem} % lang-pl
    Podzielić trójkąt prostą przez dany punkt na dwie części o równych polach. % lang-pl
    Sono dati un segmento $AB$ e un punto $P$ all'interno di un cerchio. % lang-it
    
\end{problem} % lang-pl
Hartshorne \cite[s. 44]{hartshorne2000}. % lang-pl
\begin{problem} % lang-pl
    Dany jest odcinek $AB$ oraz punkt $P$ wewnątrz okręgu. % lang-pl
    Skonstruować cięciwę tego okręgu, która przechodzi przez punkt $P$ o długości takiej samej, jak odcinek $AB$. % lang-pl
    Costruire una corda del cerchio passante per $P$ e avente la stessa lunghezza del segmento $AB$. % lang-it
    \index{cięciwa}
\end{problem}

\begin{problem}
    Dany jest odcinek $AB$, inny odcinek o długości $d$ oraz kąt $\alpha$. % lang-pl
    
    Skonstruować trójkąt $ABC$ tak, by kąt przy wierzchołku $C$ miał miarę $\alpha$, zaś suma długości ramion tego kąta była równa $d$. % lang-pl
    Sono dati un segmento $AB$, un altro segmento di lunghezza $d$ e un angolo $\alpha$. % lang-it
    
    Costruire un triangolo $ABC$ tale che l'angolo in $C$ abbia ampiezza $\alpha$ e che la somma delle lunghezze dei lati che formano tale angolo sia uguale a $d$. % lang-it
\end{problem}

\begin{problem}
    Dane są dwa okręgi takie, że żaden nie jest zawarty w drugim. % lang-pl
    
    Skonstruować styczną do obydwu okręgów. % lang-pl
    Sono dati due cerchi, nessuno dei quali è contenuto nell'altro. % lang-it
    
    Costruire una tangente comune ai due cerchi. % lang-it
    \index{dwustyczna}
\end{problem}

\begin{problem}
    Dany jest okrąg $\Gamma$ oraz jego środek $O$. % lang-pl
    
    Skonstruować trzy przystające okręgi, które są styczne do pozostałych dwóch oraz do $\Gamma$. % lang-pl
    Sono dati un cerchio $\Gamma$ e il suo centro $O$. % lang-it
    
    \hfill \emph{(13 kroków)}. % lang-pl
    Costruire tre cerchi congruenti tangenti tra loro a coppie e tangenti a $\Gamma$. % lang-it
    \hfill \emph{(13 passi)}. % lang-it
\end{problem}

\begin{problem}
    Dany jest okrąg $\Gamma$ oraz dwa punkty $A$ i $B$. % lang-pl
    
    Skonstrować punkt $C$ na okręgu $\Gamma$ tak, by odcinek łączący punkty przecięcia prostych $CA$, $CB$ z okręgiem $\Gamma$ był równoległy do odcinka $AB$. % lang-pl
    Sono dati un cerchio $\Gamma$ e due punti $A$ e $B$. % lang-it
    
    Costruire un punto $C$ sul cerchio $\Gamma$ tale che il segmento che unisce i secondi punti d'intersezione delle rette $CA$ e $CB$ con $\Gamma$ sia parallelo ad $AB$. % lang-it
\end{problem}

\begin{problem}
    Skonstruować trzy parami styczne okręgi, każdy o innym promieniu, których środki nie są współliniowe. % lang-pl
    
    \hfill \emph{(7 kroków)}. % lang-pl
    Costruire tre cerchi mutuamente tangenti, tutti di raggio diverso, con centri non allineati. % lang-it
    \hfill \emph{(7 passi)}. % lang-it
\end{problem}

Skonstruować oś potęgową. % lang-pl

Costruire l'asse radicale. % lang-it
\index{oś!potęgowa}
Audin \cite[s. 107]{audin_2003} podaje ten fakt w formie ćwiczenia. % lang-pl

Audin \cite[p. 107]{audin_2003} presenta questo fatto come esercizio. % lang-it

\begin{problem}
\label{konstrukcja_kwadratu_wpisanego}%
    Dany jest trójkąt $ABC$ z kątami ostrymi przy $B$ i $C$. % lang-pl
    
    Skonstruować kwadrat, którego jeden bok leży na odcinku $BC$, zaś pozostałe dwa wierzchołki leżą na odcinkach $AC$ i $AB$. % lang-pl
    È dato un triangolo $ABC$ con angoli acuti in $B$ e $C$. % lang-it
    
    Costruire un quadrato avente un lato sul segmento $BC$ e gli altri due vertici sui segmenti $AC$ e $AB$. % lang-it
\end{problem}

(Po tej konstrukcji warto przeczytać o trójkącie Calabiego z faktu \ref{trojkat_calabiego} raz jeszcze). % lang-pl

(Dopo questa costruzione vale la pena rileggere il fatto \ref{trojkat_calabiego} sul triangolo di Calabi). % lang-it

%